\section{MetaMem：由事实存储转向事实利用}
\label{sec:metamem}

MetaMem\upcite{xin-etal-2026-metamem}研究检索后的事实利用：模型即使找到相关记忆，仍可能误解时间、数量或从属关系。该方法将可复用的知识使用经验保存为元记忆，并通过答案反馈迭代优化。

\subsection{研究问题与两种记忆}

长期对话包含时间、人物、地点、数量及其从属关系。例如，一次旅行总里程为1200英里，其中第一天的300英里已包含在总数中；若另三次旅行合计1800英里，四次旅行应为3000英里，而非3300英里。错误不仅可能来自遗漏证据，也可能来自取回后对“整体与部分”关系的误判。

论文区分事实记忆 $M$ 与元记忆 $E$：前者保存具体经历与知识，后者保存\textbf{可迁移的知识使用经验}，例如在计数前检查是否重复包含。设问题为 $q$，主方法可抽象为
\begin{equation}
\hat y=\mathrm{LLM}\bigl(\operatorname{Prompt}(q,M,E)\bigr).
\end{equation}
普通记忆问答主要提供$(q,M)$，MetaMem进一步将$E$写入提示词。优化对象是自然语言规则，基础语言模型的权重保持不变。

\subsection{从反馈到规则的自反思优化}

方法包含规则学习与冻结后的问答两个阶段。训练时，给定带标准答案的样本$(q,a,M)$，模型在当前元记忆下多次生成回答，再利用标准答案评价。反思模块归纳成功与失败原因，提出规则的添加、删除或修改，并合并批内重复与冲突的提案。验证集用于选择规则版本，测试时使用固定规则和独立问题。

\begin{center}
\small
训练数据与候选回答 $\longrightarrow$ 正误反馈与反思 $\longrightarrow$ 规则编辑与合并
$\longrightarrow$ 冻结后的元记忆 $E$ $\longrightarrow$ 新问题的事实利用
\end{center}

规则学习的主要成本来自多次回答采样、反馈、反思与编辑；冻结后的问答阶段读取事实记忆和规则。

\subsection{实验设置与结果}

主要数据集LongMemEval包含500道问题，采用五折划分，每折350题训练、50题验证、100题测试。生成模型为Qwen3-30B-A3B-Instruct与Llama3.1-70B-Instruct，事实记忆基座主要为LightMem，默认评价模型为Qwen3-235B-A22B。规则学习默认运行5个epoch，每批50题，每题采样5次。

\begin{table}[!htbp]
\centering\small
\caption{LongMemEval平均准确率（\%）}
\begin{tabular}{lrrr}
\toprule
生成模型 & LightMem & MetaMem & 增量（百分点）\\
\midrule
Qwen3-30B-A3B-Instruct & 67.50 & 71.90 & +4.40\\
Llama3.1-70B-Instruct & 66.17 & 69.08 & +2.91\\
\bottomrule
\end{tabular}
\end{table}

两种生成模型的平均准确率分别提高4.40和2.91个百分点。Qwen配置下，完整方法为71.90\%，去除反思后为70.02\%，去除规则演化后为68.33\%，说明两个环节均有贡献。

\subsection{方法局限与课设启发}

元记忆将事实召回与事实利用分开处理。自然语言规则便于检查、合并与迁移，可用于分析时间推理、计数和跨事件引用中的重复错误。

规则学习依赖带答案样本与评价模型，错误评判可能污染更新；多次采样和编辑增加优化成本。规则过度泛化、彼此冲突或积累过长也会影响回答。方法发挥作用的前提是关键事实已进入上下文，其在小模型、小样本和其他任务上的适用性仍需检验。

StructMem侧重事实与关系的组织，MetaMem侧重检索结果的使用。课程设计可据此将失败题分为证据缺失与证据利用错误，并在固定证据的条件下检查重复计数、时间误读等问题。这样的诊断有助于判断下一步应调整检索还是回答策略。

% 原文定位：原文PDF第1页，图1；原文PDF第3页，式(1)--(3)；原文PDF第5、11页，\S4与附录A.3、A.5；原文PDF第7页，表2；原文PDF第9、11--12页及本组既有代码核查。
